\documentclass[11pt,a4paper]{article}

\usepackage{fontspec}
\setmainfont{DejaVu Sans}
\setmonofont[Scale=0.88]{DejaVu Sans Mono}

\usepackage[a4paper,margin=20mm,top=24mm,bottom=24mm,headheight=14pt]{geometry}
\usepackage{xcolor}
\usepackage{titlesec}
\usepackage{fancyhdr}
\usepackage{listings}
\usepackage{tcolorbox}
\usepackage{booktabs}
\usepackage{tabularx}
\usepackage{hyperref}
\usepackage{microtype}

% --- Кольорова палітра ---
\definecolor{primary}{HTML}{0969DA}
\definecolor{primarydark}{HTML}{0B5CAD}
\definecolor{ink}{HTML}{1F2328}
\definecolor{muted}{HTML}{57606A}
\definecolor{codebg}{HTML}{F6F8FA}
\definecolor{codeframe}{HTML}{D0D7DE}
\definecolor{keywordcolor}{HTML}{CF222E}
\definecolor{stringcolor}{HTML}{0A3069}
\definecolor{commentcolor}{HTML}{6E7781}
\definecolor{tipborder}{HTML}{2DA44E}
\definecolor{tipbg}{HTML}{F0FFF4}
\definecolor{noteborder}{HTML}{0969DA}
\definecolor{notebg}{HTML}{F0F8FF}

% --- Налаштування заголовків ---
\titleformat{\section}{\Large\bfseries\color{primarydark}}{§\,\thesection.}{0.6em}{}[\vspace{1pt}\color{codeframe}\hrule\vspace{4pt}]
\titleformat{\subsection}{\large\bfseries\color{ink}}{\thesubsection}{0.6em}{}
\titleformat{\subsubsection}{\normalsize\bfseries\color{primarydark}}{\thesubsubsection}{0.5em}{}

% --- Верхній і нижній колонтитули ---
\pagestyle{fancy}
\fancyhf{}
\fancyhead[L]{\small\color{muted}Free Pascal: Посібник із сучасної конкурентності}
\fancyhead[R]{\small\color{muted}Гілка \texttt{feature/concurrency}}
\fancyfoot[C]{\small\thepage}
\renewcommand{\headrulewidth}{0.4pt}
\renewcommand{\footrulewidth}{0.4pt}

% --- Налаштування підсвічування коду ---
\lstdefinelanguage{ModernPascal}{
  morekeywords={
    program,unit,interface,implementation,uses,begin,end,type,var,const,
    generic,specialize,function,procedure,class,record,object,
    override,overload,sealed,helper,for,in,to,downto,do,while,repeat,until,
    if,then,else,case,of,try,except,finally,raise,
    async,await,reference,to,out,inline,constructor,destructor,property,read,write
  },
  morekeywords=[2]{
    ITask,IFuture,TTask,TParallel,TCancellationToken,TCancellationTokenSource,
    TChannel,TSlimReaderWriterLock,TCountdownEvent,TBarrier,
    IEventLoop,TEventLoop,TEventLoopGroup,IAsyncSocket,IAsyncServerSocket,
    TFiber,TGenerator,TTaskCompletionSource,TThreadPool
  },
  morekeywords=[3]{
    implicitgenerics,anonymousfunctions,functionreferences,advancedrecords
  },
  sensitive=false,
  morecomment=[l]{//},
  morecomment=[s]{/*}{*/},
  morecomment=[s]{\{}{\}},
  morestring=[b]',
}

\lstset{
  language=ModernPascal,
  backgroundcolor=\color{codebg},
  frame=single,
  rulecolor=\color{codeframe},
  basicstyle=\ttfamily\small\color{ink},
  keywordstyle=\bfseries\color{keywordcolor},
  keywordstyle=[2]\bfseries\color{primarydark},
  keywordstyle=[3]\color{primary},
  commentstyle=\itshape\color{commentcolor},
  stringstyle=\color{stringcolor},
  numbers=left,
  numberstyle=\tiny\color{muted},
  numbersep=7pt,
  breaklines=true,
  breakatwhitespace=true,
  keepspaces=true,
  columns=flexible,
  tabsize=2,
  showstringspaces=false,
  xleftmargin=12pt,
  framexleftmargin=8pt,
  aboveskip=8pt,
  belowskip=8pt
}

% --- Блоки приміток ---
\newtcolorbox{infobox}[2][]{
  colback=notebg,
  colframe=noteborder,
  coltitle=white,
  fonttitle=\bfseries,
  title={#2},
  boxrule=0.8pt,
  arc=3pt,
  left=8pt,right=8pt,top=6pt,bottom=6pt,
  #1
}

\newtcolorbox{tipbox}[2][]{
  colback=tipbg,
  colframe=tipborder,
  coltitle=white,
  fonttitle=\bfseries,
  title={#2},
  boxrule=0.8pt,
  arc=3pt,
  left=8pt,right=8pt,top=6pt,bottom=6pt,
  #1
}

% --- Налаштування Hyperref ---
\hypersetup{
  colorlinks=true,
  linkcolor=primarydark,
  urlcolor=primary,
  citecolor=primarydark,
  pdftitle={Free Pascal: Довідник і посібник з конкурентності, асинхронного I/O та узагальнень},
  pdfauthor={Free Pascal Core Team / feature/concurrency},
  pdfsubject={Concurrency, Async I/O, Generics, EventLoop, Fibers, CSP Channels}
}

\begin{document}

% --- ТИТУЛЬНА СТОРІНКА ---
\begin{titlepage}
  \centering
  \vspace*{15mm}
  {\Huge\bfseries\color{primarydark} Free Pascal\par}
  \vspace{4mm}
  {\LARGE\bfseries Сучасна модель конкурентності, асинхронного I/O та узагальнень\par}
  \vspace{2mm}
  {\large Професійний посібник для користувача-розробника з описом нововведень\par}
  \vspace{10mm}
  \rule{0.6\textwidth}{1.2pt}\par
  \vspace{10mm}
  
  \begin{tabular}{ll}
    \textbf{Цільова версія:} & Free Pascal 3.3.1 (Trunk) \\
    \textbf{Робоча гілка:} & \texttt{feature/concurrency} \\
    \textbf{Охоплені етапи:} & Фази 1.1, 1.2, 1.3, 1.4 та Розділ 2.1 \\
    \textbf{Підтримувані ОС:} & Linux (\texttt{io\_uring}, \texttt{epoll}), macOS / BSD (\texttt{kqueue}), Windows (\texttt{IOCP}) \\
    \textbf{Архітектура CPU:} & x86\_64 (AMD64 System V ABI / Win64) та AArch64 \\
  \end{tabular}
  
  \vfill
  
  \begin{infobox}{Ключові архітектурні зміни гілки}
    \small
    \textbf{1. Синтаксис мови:} \texttt{\{\$modeswitch implicitgenerics\}} дозволяє повністю відкинути архаїчне ключове слово \texttt{specialize} та вмикає автоматичне виведення типів для узагальнених методів і анонімних функцій.\\
    \textbf{2. Оркестрація задач:} \texttt{System.Threading} з неблокувальними продовженнями (\texttt{.ThenBy}, \texttt{.ContinueWith}) та токенами кооперативного скасування (\texttt{TCancellationToken}).\\
    \textbf{3. Потокобезпечна комунікація:} Модель каналів CSP (\texttt{TChannel<T>}) у стилі мови Go та ультрашвидкий \texttt{TSlimReaderWriterLock} із RAII-гардами.\\
    \textbf{4. Реактивний I/O:} Пакет \texttt{fcl-eventloop} на нативних мультиплексорах ядра ОС для обслуговування 100\,000+ з'єднань.\\
    \textbf{5. Корутини:} Пакет \texttt{fcl-async} (\texttt{TFiber}, генератори) та підтримка \texttt{async}/\texttt{await} на рівні компілятора.
  \end{infobox}
  
  \vspace{10mm}
  {\small\color{muted} 2026 рік}
\end{titlepage}

\tableofcontents
\newpage

% --- РОЗДІЛ 1 ---
\section{Архітектурна модель та вибір парадигми}

Сучасний Free Pascal у гілці \texttt{feature/concurrency} надає розробнику цілісний трирівневий набір інструментів, що усуває потребу в прямій низькорівневій роботі з системними потоками (\texttt{TThread}) та блокуючими м'ютексами.

\begin{table}[h!]
\centering
\small
\begin{tabularx}{\textwidth}{p{3.2cm}p{3.5cm}p{4.2cm}X}
\toprule
\textbf{Клас задачі} & \textbf{Основний модуль} & \textbf{Ключові класи/типи} & \textbf{Переваги} \\
\midrule
\textbf{CPU-Bound} (обчислення, стиснення, парсинг) & \texttt{System.Threading} & \texttt{TParallel.For}, \texttt{ITask}, \texttt{IFuture<T>}, \texttt{ThenBy} & Автоматичне балансування навантаження (\emph{work stealing}), мінімум блокувань. \\
\addlinespace
\textbf{Синхронізація} (обмін повідомленнями, буферизація) & \texttt{SyncObjs} & \texttt{TChannel<T>}, \texttt{TSlimReaderWriterLock}, \texttt{TBarrier} & CSP-комунікація, усунення гонок даних, відсутність дескрипторів ядра ОС. \\
\addlinespace
\textbf{I/O-Bound} (High-Load мережа, сокети, таймери) & \texttt{fcl-eventloop} & \texttt{IEventLoop}, \texttt{TEventLoopGroup}, \texttt{IAsyncSocket} & Нативні мультиплексори (\texttt{io\_uring}, \texttt{epoll}), 100\,000+ з'єднань на групу потоків. \\
\addlinespace
\textbf{Корутини} (лінійний код для асинхронних операцій) & \texttt{fcl-async} та компілятор & \texttt{TFiber}, \texttt{TGenerator<T>}, \texttt{async} / \texttt{await} & Стек файбера перемикається за $<$10\,нс без переривання системного потоку ОС. \\
\bottomrule
\end{tabularx}
\caption{Матриця підбору парадигми конкурентності}
\end{table}

\newpage

% --- РОЗДІЛ 2 ---
\section{Синтаксична ергономіка: \texttt{\{\$modeswitch implicitgenerics\}}}

\subsection{Опущення ключового слова \texttt{specialize}}
Раніше в режимі \texttt{\{\$mode objfpc\}} використання будь-якого узагальненого типу або виклику методу вимагало громіздкої конструкції \texttt{specialize}. 

Перемикач \texttt{\{\$modeswitch implicitgenerics\}} дозволяє використовувати звичний, лаконічний синтаксис для **всіх** типів (як нових, так і стандартних FCL):

\begin{lstlisting}
{$mode objfpc}{$H+}
{$modeswitch implicitgenerics}

uses fgl, Generics.Collections, System.Threading;

type
  { Оголошення типів без specialize }
  TIntList   = TFPGList<Integer>;
  TUserCache = TDictionary<string, Int64>;

var
  List: TFPGList<Integer>;
  Map:  TDictionary<string, Int64>;
  Fut:  IFuture<string>;
begin
  List := TFPGList<Integer>.Create;
  Map  := TDictionary<string, Int64>.Create;
  
  { Явний виклик generic-методу без specialize }
  Fut := TTask.Future<string>(
    function: string
    begin
      Result := 'Дані готові';
    end);
end.
\end{lstlisting}

\subsection{Автоматичне виведення типів (Generic Type Inference)}
Компілятор автоматично визначає параметри \texttt{<T>} із переданих фактичних аргументів. Це стосується чисел, рядків, інтерфейсів, об'єктів та, що найважливіше, **типів повернення анонімних функцій**:

\begin{lstlisting}
var
  FutStr: IFuture<string>;
  FutInt: IFuture<Integer>;
begin
  // T = string виводиться автоматично з типу повернення:
  FutStr := TTask.Future(
    function: string
    begin
      Result := 'Привіт, світ';
    end);

  // T = Integer виводиться автоматично:
  FutInt := TTask.Future(
    function: Integer
    begin
      Result := 42;
    end);
end.
\end{lstlisting}

\begin{tipbox}{Працює для всіх функцій і методів}
Type Inference не обмежується лише тасками. Будь-яка узагальнена функція на зразок \texttt{generic function MaxVal<T>(A, B: T): T} тепер викликається просто: \texttt{M := MaxVal(3.14, 2.71);} без вказування \texttt{<Double>}.
\end{tipbox}

\newpage

% --- РОЗДІЛ 3 ---
\section{Задачі, конвеєри та скасування (\texttt{System.Threading})}

\subsection{Кооперативне скасування: \texttt{TCancellationTokenSource}}
Джерело токенів скасування (\texttt{TCancellationTokenSource}) генерує легковагові структури \texttt{TCancellationToken}, які передаються в операції для їх узгодженого переривання:

\begin{lstlisting}
var
  CTS: TCancellationTokenSource;
  Token: TCancellationToken;
  Task: ITask;
begin
  CTS := TCancellationTokenSource.Create;
  Token := CTS.Token;

  Task := TTask.Run(
    procedure
    begin
      while not Token.IsCancellationRequested do
      begin
        DoWorkBatch;
      end;
      Token.ThrowIfCancellationRequested; // Генерує EOperationCancelled
    end, Token);

  // Скасування операції з іншого потоку або за таймаутом:
  CTS.CancelAfter(500); // скасувати за 500 мс
end.
\end{lstlisting}

\subsection{Неблокувальні ланцюжки задач: \texttt{.ThenBy} та \texttt{.ContinueWith}}
Замість блокування потоку у виклику \texttt{.Wait} або \texttt{.Value}, використовується побудова неблокувальних конвеєрів:

\begin{lstlisting}
var
  DownloadTask: IFuture<string>;
  ParseTask:    IFuture<Integer>;
begin
  // 1. Початкова задача
  DownloadTask := TTask.Future(
    function: string
    begin
      Result := HttpGet('https://api.service.internal/count');
    end);

  // 2. Неблокувальне продовження ThenBy (TResult = Integer виведено автоматично)
  ParseTask := DownloadTask.ThenBy(
    function(const Prev: ITask): Integer
    begin
      Result := StrToIntDef(IFuture<string>(Prev).Value, 0);
    end);

  // 3. Фінальна обробка результату або винятку
  ParseTask.ContinueWith(
    procedure(const Completed: ITask)
    begin
      if Completed.Status = TTaskStatus.Completed then
        WriteLn('Кількість: ', IFuture<Integer>(Completed).Value)
      else
        WriteLn('Помилка: ', Completed.ExceptionObject.Message);
    end);
end.
\end{lstlisting}

\newpage

% --- РОЗДІЛ 4 ---
\section{Нові примітиви синхронізації (\texttt{SyncObjs})}

\subsection{Канали CSP: \texttt{generic TChannel<T>}}
Канал реалізує модель Communicating Sequential Processes (як у мові Go). Він підтримує як обмежену буферизацію (\texttt{Capacity > 0} з автоматичним блокуванням записувачів при заповненні), так і необмежену.

\begin{lstlisting}
var
  Channel: TChannel<string>;
begin
  Channel := TChannel<string>.Create(50); // Буфер на 50 повідомлень

  // Письменник:
  TTask.Run(
    procedure
    var I: Integer;
    begin
      for I := 1 to 1000 do
        Channel.Send('Повідомлення #' + IntToStr(I)); // Backpressure
      Channel.Close;
    end);

  // Читач з прямим обходом for..in без накладних алокацій:
  TTask.Run(
    procedure
    var Item: string;
    begin
      for Item in Channel do
      begin
        ProcessItem(Item);
      end;
      WriteLn('Канал закрито, обробку завершено.');
    end);
end.
\end{lstlisting}

\subsection{\texttt{TSlimReaderWriterLock} та RAII-гарди}
Надвисокопродуктивне блокування читання/запису без системних дескрипторів ядра ОС. Підтримує пріоритет записувачів, щоб запобігти голодуванню, та RAII-гарди:

\begin{lstlisting}
var
  Lock: TSlimReaderWriterLock;

procedure ReadResource;
var
  Guard: TSlimReaderWriterLock.IReadLock;
begin
  Guard := Lock.ReadLock; // Захоплює дозвіл на читання
  // Ресурс захищено. Звільнення відбувається автоматично при виході зі скоупу!
end;

procedure WriteResource;
var
  Guard: TSlimReaderWriterLock.IWriteLock;
begin
  Guard := Lock.WriteLock; // Ексклюзивний дозвіл на запис
  UpdateResourceData;
end;
\end{lstlisting}

\newpage

% --- РОЗДІЛ 5 ---
\section{Масштабований асинхронний I/O (\texttt{fcl-eventloop})}

Пакет \texttt{fcl-eventloop} реалізує патерн Reactor поверх найбільш ефективних нативних API операційної системи:
\begin{itemize}
  \item \textbf{Linux:} \texttt{io\_uring} (для ядер 5.1+) та високошвидкісний \texttt{epoll};
  \item \textbf{macOS / FreeBSD:} \texttt{kqueue};
  \item \textbf{Windows:} \texttt{IOCP} (I/O Completion Ports).
\end{itemize}

\subsection{Побудова сервера на базі \texttt{IEventLoop}}
\begin{lstlisting}
var
  Loop:   IEventLoop;
  Server: IAsyncServerSocket;
begin
  Loop := TEventLoop.Create(bkAuto); // Автоматичний вибір кращого бекенду
  Server := TEventLoop.CreateTcpServer(Loop);
  Server.Listen('0.0.0.0', 8080);

  Server.AsyncAccept(
    procedure(AHandle: TLoopHandle; const Addr: string; Port: Word; Status: Integer)
    var
      Client: IAsyncSocket;
      Buffer: array[0..4095] of Byte;
    begin
      if Status <> 0 then Exit;
      Client := TAsyncSocket.Create(Loop, AHandle, True);
      
      Client.AsyncRead(@Buffer[0], SizeOf(Buffer),
        procedure(Bytes: NativeInt; ReadStatus: Integer)
        begin
          if (ReadStatus = 0) and (Bytes > 0) then
          begin
            // Неблокувальне відлуння (Echo)
            Client.AsyncWrite(@Buffer[0], Bytes, nil);
          end
          else
            Client.Close;
        end);
    end);

  Loop.Run;
end.
\end{lstlisting}

\subsection{\texttt{TEventLoopGroup}: Multi-Reactor на 100\,000+ з'єднань}
Клас \texttt{TEventLoopGroup} створює пул незалежних потоків-реакторів (наприклад, по одному на кожне апаратне ядро процесора). Нові з'єднання балансуються між реакторами методом Round-Robin або Least-Connections, усуваючи конкуренцію за чергу подій.

\newpage

% --- РОЗДІЛ 6 ---
\section{Корутини, файбери та \texttt{async}/\texttt{await} (\texttt{fcl-async})}

\subsection{Легковагові стекові файбери: \texttt{TFiber}}
\texttt{TFiber} дозволяє писати лінійний код, який призупиняється в очікуванні асинхронних операцій без блокування системного потоку ОС.
\begin{itemize}
  \item Контекст перемикається на вбудованому асемблері AMD64 менш ніж за 10\,нс;
  \item Гарантується 16-байтне вирівнювання стеку під SIMD/AVX інструкції;
  \item Стек винятків FPC (\texttt{ExceptAddrStack}) симетрично ізолюється для кожного файбера.
\end{itemize}

\begin{lstlisting}
var
  WorkerFiber: TFiber;
begin
  WorkerFiber := TFiber.Create(
    procedure
    var
      Fut: IFuture<string>;
      Data: string;
    begin
      WriteLn('Файбер розпочав виконання');

      Fut := TTask.Future(
        function: string
        begin
          Sleep(200);
          Result := 'Завантажені дані';
        end);

      // Призупиняє ВИКЛЮЧНО поточний файбер, потік ОС продовжує іншу роботу!
      Data := TFiber.Await(Fut);

      WriteLn('Файбер відновився з результатом: ', Data);
    end);

  WorkerFiber.Start;
end.
\end{lstlisting}

\subsection{Ледачі послідовності: \texttt{generic TGenerator<T>}}
Генератори реалізують патерн ітератора, де кожне наступне значення обчислюється лише тоді, коли його запитує цикл \texttt{for..in}:

\begin{lstlisting}
function GenerateFibonacci(Limit: Integer): TGenerator<Int64>;
begin
  Result := TGenerator<Int64>.Create(
    procedure(const Yield: TGenerator<Int64>.TYieldProc)
    var A, B, C: Int64;
    begin
      A := 0; B := 1;
      while A <= Limit do
      begin
        Yield(A); // Призупиняє генератор і повертає значення у for..in
        C := A + B; A := B; B := C;
      end;
    end);
end;

var Num: Int64;
begin
  for Num in GenerateFibonacci(1000) do
    WriteLn(Num);
end.
\end{lstlisting}

\newpage

% --- РОЗДІЛ 7 ---
\section{Комплексний наскрізний приклад}

Наведений нижче приклад демонструє спільну роботу всіх компонентів: високонавантажений TCP-сервер на \texttt{fcl-eventloop} приймає байти, через \texttt{TChannel<string>} передає їх у пул потоків \texttt{System.Threading}, де вони паралельно трансформуються через ланцюжок \texttt{.ThenBy} з автоматичним виведенням типів:

\begin{lstlisting}
program EndToEndModernConcurrency;
{$mode objfpc}{$H+}
{$modeswitch implicitgenerics}
{$modeswitch anonymousfunctions}
{$modeswitch functionreferences}

uses
  {$IFDEF UNIX}cthreads,{$ENDIF}
  SysUtils, SyncObjs, System.Threading,
  eventloop.types, eventloop.core, eventloop.sockets, fcleventloop;

var
  WorkQueue: TChannel<string>;
  Loop:      IEventLoop;
  Server:    IAsyncServerSocket;
  CTS:       TCancellationTokenSource;

begin
  CTS := TCancellationTokenSource.Create;
  WorkQueue := TChannel<string>.Create(1000); // Буфер із backpressure

  // 1. Робочий потік-диспетчер пулу задач:
  TTask.Run(
    procedure
    var Msg: string;
    begin
      for Msg in WorkQueue do
      begin
        // Обчислення в пулі потоків з автоматичним виведенням типів
        TTask.Future(
          function: Integer
          begin
            Result := Length(Msg);
          end
        ).ThenBy(
          function(const Prev: ITask): string
          begin
            Result := 'Довжина вхідного пакета: ' + IntToStr(IFuture<Integer>(Prev).Value);
          end
        ).ContinueWith(
          procedure(const T: ITask)
          begin
            WriteLn(IFuture<string>(T).Value);
          end);
      end;
    end, CTS.Token);

  // 2. Реактор мережевого I/O (fcl-eventloop):
  Loop := TEventLoop.Create(bkAuto);
  Server := TEventLoop.CreateTcpServer(Loop);
  Server.Listen('0.0.0.0', 9000);

  Server.AsyncAccept(
    procedure(AHandle: TLoopHandle; const Addr: string; Port: Word; Status: Integer)
    var
      Sock: IAsyncSocket;
      Buf: array[0..2047] of AnsiChar;
    begin
      if Status <> 0 then Exit;
      Sock := TAsyncSocket.Create(Loop, AHandle, True);
      Sock.AsyncRead(@Buf[0], SizeOf(Buf),
        procedure(Bytes: NativeInt; ReadStatus: Integer)
        begin
          if (ReadStatus = 0) and (Bytes > 0) then
            WorkQueue.Send(Copy(Buf, 1, Bytes))
          else
            Sock.Close;
        end);
    end);

  WriteLn('Сервер запущено на порті 9000. Очікування клієнтів...');
  Loop.Run;

  WorkQueue.Close;
  CTS.Cancel;
  CTS.Free;
  WorkQueue.Free;
end.
\end{lstlisting}

\newpage

% --- РОЗДІЛ 8 ---
\section{Правила надійності та перевірка під Valgrind}

\begin{enumerate}
  \item \textbf{Ніколи не блокуйте потік реактора \texttt{IEventLoop}:} Будь-які блокуючі системні виклики (робота з диском, звернення до СУБД, важкі обчислення) повинні делегуватися у пул потоків через \texttt{TTask.Run} або передаватися в чергу \texttt{TChannel<T>}. Потік \texttt{EventLoop} повинен лише мультиплексувати дескриптори.
  \item \textbf{Ініціалізація \texttt{cthreads} на Linux/Unix:} Будь-яка багатопотокова програма зобов'язана першим модулем у секції \texttt{uses} підключати модуль \texttt{cthreads}, інакше внутрішні примітиви синхронізації RTL працюватимуть некоректно.
  \item \textbf{Перевірка детекторами гонок даних (Helgrind / DRD):} Усі нові примітиви (\texttt{TCancellationToken}, \texttt{TSlimReaderWriterLock}, \texttt{TChannel}, \texttt{TCountdownEvent}, \texttt{TBarrier}) підтримують ключ компілятора \texttt{-dFPC\_VALGRIND}. У бінарний код вбудовуються клієнтські запити Valgrind (\emph{happens-before / happens-after}), що гарантує повну відсутність хибних спрацювань під час аналізу.
\end{enumerate}

\begin{infobox}{Корисні посилання в репозиторії}
\begin{itemize}
  \item \texttt{docs/concurrency/modern-concurrency-reference.md} --- повний довідник API;
  \item \texttt{docs/concurrency/guide-async.md} --- файбери та корутини;
  \item \texttt{docs/concurrency/guide-eventloop.md} --- цикл подій та архітектура реактора;
  \item \texttt{docs/concurrency/roadmap-improvements.md} --- дорожня карта гілки \texttt{feature/concurrency}.
\end{itemize}
\end{infobox}

\end{document}
